\section{Dense meridians with no quadrilateral type guesses}
\label{sec:fibonacci}

The support obstruction in the previous continuation does not obstruct
mandatory-type propagation.  We prove this for the entire Fibonacci layered
solid-torus family, starting with all three quadrilateral types available in
every tetrahedron.  The proof supplies exact identities that certify the
propagation; it does not extrapolate from a finite collection of LP runs.
Throughout this section, $\beta_{\mathrm{LP}}$ refers to exhaustive
\emph{mandatory-type} propagation, without an additional rule that explicitly
deletes every coordinate forced to zero by the relaxed cone.

\subsection{The labelled triangulations and their full relaxation}

For $t\ge1$, number the tetrahedra $0,\ldots,t-1$ and their local vertices
$0,1,2,3$.  A permutation below lists the images of these four vertices;
in particular, it determines the target face of a gluing.  For each
$0\le i<t-1$, make the identifications
\begin{equation}
\begin{aligned}
 (i,\text{face }0)&\longrightarrow
       (i+1,\text{face }2),&\qquad &(2,1,3,0),\\
 (i,\text{face }1)&\longrightarrow
       (i+1,\text{face }3),& &(0,3,1,2).
\end{aligned}
\label{eq:fib-face-gluings}
\end{equation}
In the final tetrahedron, glue face $0$ to face $1$ using
$(1,2,3,0)$.  Include the inverse face maps, and leave faces $2,3$ of
tetrahedron zero unpaired.  Denote the resulting triangulation by $B_t$.
This is the same explicit family as the previous package's
\texttt{fibonacci\_fixture.py}: it is a one-vertex triangulation of a solid
torus, obtained by repeated standard layerings on the one-tetrahedron solid
torus~\cite{priorKernel}.  Its two unpaired faces form the boundary torus.
The face maps, rather than an isomorphism signature or a numerical size
parameter, specify the family for arbitrary $t$.

Put $\ell=t-1$.  Write the four triangle coordinates in tetrahedron $i$ as
$(A_i,B_i,C_i,D_i)$, and its quadrilateral coordinates, in the implementation's
type order $0,1,2$, as $(U_i,V_i,W_i)$.  At this point all these variables are
merely nonnegative real numbers.  We impose the standard matching equations
but allow incompatible quadrilateral types to coexist.  The Euler functional
$\chi$ on this relaxation means the usual linear extension of the normal
surface Euler characteristic; a relaxed vector need not yet represent an
embedded surface.

\begin{lemma}[Full relaxed matching equations]
\label{lem:fib-relaxed}
Every relaxed matching vector on $B_t$ has
\[
 A_i=B_i=X_i,\qquad C_i=D_i=Y_i.
\]
In these variables, its matching equations are exactly
\begin{equation}
\begin{aligned}
 X_i+U_i&=X_{i+1}+W_{i+1},\\
 Y_i+V_i&=Y_{i+1}+U_{i+1},\\
 Y_i+W_i&=X_{i+1}+V_{i+1}
       &&(0\le i<\ell),\\
 X_\ell+U_\ell&=Y_\ell+W_\ell.
\end{aligned}
\label{eq:fib-relaxed}
\end{equation}
\end{lemma}

\begin{proof}
The six face-corner equations between tetrahedra $i$ and $i+1$ are
\[
\begin{array}{rcl@{\qquad}rcl}
 B_i+U_i&=&B_{i+1}+W_{i+1},&
 A_i+U_i&=&A_{i+1}+W_{i+1},\\
 C_i+V_i&=&D_{i+1}+U_{i+1},&
 D_i+V_i&=&C_{i+1}+U_{i+1},\\
 D_i+W_i&=&A_{i+1}+V_{i+1},&
 C_i+W_i&=&B_{i+1}+V_{i+1}.
\end{array}
\]
The final self-gluing gives
\[
 B_\ell+U_\ell=C_\ell+W_\ell,\qquad
 C_\ell=D_\ell,\qquad
 A_\ell+U_\ell=D_\ell+W_\ell.
\]
Consequently $A_\ell=B_\ell$ and $C_\ell=D_\ell$.  Subtracting the two
equations in each of the first two rows of the array propagates these
equalities backwards through all tetrahedra.  The remaining equations reduce
to \eqref{eq:fib-relaxed}.  Conversely, the displayed pair equalities and
\eqref{eq:fib-relaxed} imply every original face-corner equation.  Thus this
is an exact description of the full relaxation, with no quadrilateral type
already selected.
\end{proof}

\subsection{Euler identities that certify mandatory types}

\begin{lemma}[Euler functional]
\label{lem:fib-euler}
On the entire relaxed matching cone,
\begin{equation}
 \chi=X_\ell-\sum_{j=0}^{\ell-1}U_j.
 \label{eq:fib-euler}
\end{equation}
An empty sum is zero, so this also covers $t=1$.
\end{lemma}

\begin{proof}
First consider the final self-glued tetrahedron by itself.  Three global
edge representatives are $01,02,03$.  With $S=U+V+W$, their total normal
intersection weight is $4X+2Y+2S$.  The paired face $0$ and the two boundary
faces $2,3$ have total arc count $5X+4Y+3S$.  The normal-disc count is
$2X+2Y+S$.  The alternating sum of normal vertices, edges and faces is
therefore $X$.

Next add an outer tetrahedron $i$ along its faces $0,1$.  The old global
edge classes persist, and the only new edge is $01$, which belongs to
neither attaching face.  The common attaching edge $23$ maps to the already
identified inner edges $03$ and $12$; hence the attachment does not merge
two previously distinct old edge classes.  The new edge contributes
$2X_i+V_i+W_i$ normal vertices.  The two old boundary faces become paired
faces and continue to be counted once, with their arc counts unchanged by
matching.  The two newly exposed faces $2,3$ each contribute
$2X_i+Y_i+U_i+V_i+W_i$ normal edges.  Finally, the new tetrahedron contributes
$2X_i+2Y_i+U_i+V_i+W_i$ normal discs.  Its net Euler contribution is
\[
\begin{aligned}
 &(2X_i+V_i+W_i)
 -2(2X_i+Y_i+U_i+V_i+W_i)\\
 &\hspace{35mm}+(2X_i+2Y_i+U_i+V_i+W_i)=-U_i.
\end{aligned}
\]
Induction proves \eqref{eq:fib-euler}.  All these are linear cell-count
identities, so they remain valid for the relaxation.  The argument does
not assume that the two boundary triangles being attached form an embedded
disc: their additional edge identifications have been accounted for in the
global edge count.
\end{proof}

Telescoping the first equation of \eqref{eq:fib-relaxed} now yields
\begin{equation}
 X_i=\chi+\sum_{j<i}U_j+\sum_{j>i}W_j\ge\chi
 \qquad(0\le i\le\ell).
 \label{eq:fib-X-dual}
\end{equation}
For $i<\ell$, the third matching equation gives the second identity
\begin{equation}
 \chi=Y_i+W_i-V_{i+1}
       -\sum_{j\le i}U_j-\sum_{j\ge i+2}W_j
 \le Y_i+W_i.
 \label{eq:fib-Y-dual}
\end{equation}
At the last tetrahedron its counterpart is
\begin{equation}
 \chi=Y_\ell+W_\ell-\sum_{j=0}^{\ell}U_j
 \le Y_\ell+W_\ell.
 \label{eq:fib-last-dual}
\end{equation}
These identities are symbolic Farkas certificates.  Each is obtained by a
linear combination of matching equalities and the Euler identity, with the
remaining terms explicitly nonnegative.  They therefore justify exact
infeasibility tests for every $t$, independently of floating-point LP output.

\begin{theorem}[Zero propagation backdoor]
\label{thm:fib-zero-backdoor}
For every $t\ge1$,
\[
 \beta_{\mathrm{LP}}(B_t)=0.
\]
More explicitly, for every choice of a zero triangle anchor at the unique
global vertex, exhaustive mandatory-type propagation on the full three-type
relaxation either proves that positive Euler characteristic is impossible
or reduces the allowed quadrilateral types to type $2$ in every tetrahedron.
No quadrilateral type need be guessed.
\end{theorem}

\begin{proof}
Every original triangle anchor is, by Lemma~\ref{lem:fib-relaxed}, an
equation $X_i=0$ or $Y_i=0$.  If $X_i=0$, equation
\eqref{eq:fib-X-dual} proves $\chi\le0$, so the positive-Euler branch is
infeasible immediately.

Suppose instead that the anchor is $Y_i=0$.  Equations
\eqref{eq:fib-Y-dual} and \eqref{eq:fib-last-dual} show that every relaxed
point with $\chi>0$ has $W_i\ge\chi>0$.  Thus $W_i$ is mandatory.
Equivalently, adding $W_i=0$ makes the positive-Euler LP infeasible, with
the displayed identity supplying its dual justification.  An admissible
surface must consequently have $U_i=V_i=0$.

These two type deletions force neighbouring triangle values to vanish.
Indeed, whenever the indicated neighbours exist, the middle matching
equation of \eqref{eq:fib-relaxed} gives
\[
 Y_{i-1}+V_{i-1}=Y_i+U_i=0,
 \qquad
 Y_{i+1}+U_{i+1}=Y_i+V_i=0.
\]
Nonnegativity implies $Y_{i-1}=Y_{i+1}=0$.  At either newly reached
tetrahedron, the same Euler inequalities make its $W$ coordinate mandatory,
and its $U,V$ coordinates are discarded.  Inducting along the chain reaches
every tetrahedron after at most $t$ propagation rounds.

The resulting allowed mask contains only $W_i$ at each tetrahedron.  This
argument uses only the mandatory-type rule.  The neighbouring triangle
zeros need not be installed as additional implementation rules: they are
already consequences of the matching equations and previous type deletions,
and every subsequent LP test enforces those equations.  Hence a fair scan
of the available mandatory-type tests detects the same implications.
All anchors have been covered, so the empty set is a uniform strong
backdoor, proving the asserted value of $\beta_{\mathrm{LP}}$.
\end{proof}

The proof also specifies a particularly simple order for a generic search.
In the original coordinate order, the first two anchors are $A_0=0$ and
$B_0=0$; both are rejected by \eqref{eq:fib-X-dual}.  The next anchor is
$C_0=Y_0=0$.  Starting there, the mandatory implications can be taken in
the order $W_0,W_1,\ldots,W_\ell$.  At each step the current matching
equations imply the next zero triangle, and the next mandatory-type test
has an infeasibility proof from \eqref{eq:fib-Y-dual} or
\eqref{eq:fib-last-dual}.  The search is not supplied the Fibonacci
coordinates or the final type sequence; these are deductions from its
unrestricted input equations.

A straightforward implementation that scans all remaining quadrilateral
coordinates after every successful implication makes at most $O(t^2)$
LP calls on this positive anchor.  There are at most $t$ successful type
choices and at most $3t$ tests per scan.  The preceding two rejected
anchors add only two feasibility calls.  With a polynomial-time rational
LP oracle this is a polynomial bit-time discovery procedure on the
family.  This statement concerns the algorithmic oracle bound; it does
not give a polynomial worst-case pivot bound for a particular simplex
implementation.  An implementation may also terminate earlier when its
current relaxed basic solution already satisfies the quadrilateral
constraints, without computing the whole closure.

\subsection{The reconstructed essential disc}

After a positive branch has propagated, all $Y_i,U_i,V_i$ vanish.  The
remaining equations are
\[
 X_i=X_{i+1}+W_{i+1},\qquad W_i=X_{i+1}\quad(i<\ell),
 \qquad X_\ell=W_\ell.
\]
Write $F_0=0$, $F_1=F_2=1$ and $F_{m+2}=F_{m+1}+F_m$.  Setting
$u=W_\ell$ and solving backwards gives
\begin{equation}
 X_i=F_{t-i+1}u,\qquad W_i=F_{t-i}u,\qquad u\ge0.
 \label{eq:fib-recovered-ray}
\end{equation}
The surviving compatible cone is therefore one-dimensional.  Every
canonical admissible positive-Euler vector belongs to it, since it has
some zero triangle anchor and every positive anchored branch forces the
same vanishings.  Its primitive integral vector is obtained with $u=1$,
since its final $W$ coordinate is
itself $u$.  In the seven-coordinate convention it is
\begin{equation}
 x_i=(F_{t-i+1},F_{t-i+1},0,0;0,0,F_{t-i}),
 \qquad 0\le i<t.
 \label{eq:fib-standard-vector}
\end{equation}

\begin{proposition}[The recovered ray is an essential disc]
\label{prop:fib-essential-disc}
The vector in \eqref{eq:fib-standard-vector} represents a connected normal
disc whose boundary is essential on $\partial B_t$.
\end{proposition}

\begin{proof}
It spans a ray of a one-dimensional coordinate face of the standard
matching cone, so it is standard extreme.  Its primitive integral vector
is connected: a decomposition into component vectors would express it as
a sum of positive integer multiples of itself.  Here integrality of the
multiples follows from a B\'ezout combination of its primitive coordinates,
and two positive integers cannot sum to one.  This argument does not assume
that a normal surface is two-sided or orientable.

Equation~\eqref{eq:fib-euler} gives $\chi=1$.  The exposed faces have
nonzero normal arcs, so the surface has boundary.  Classification of compact
connected surfaces therefore makes it a disc.  Its three boundary-edge
intersection weights, in some order, are
\[
 F_{t+1},\qquad F_{t+2},\qquad F_{t+3}.
\]
For example, at the outer tetrahedron the weights on $03$, $02$ and $01$
are respectively $X_0$, $X_0+W_0$ and $2X_0+W_0$, giving these numbers.
Consecutive Fibonacci numbers are coprime, so these weights are not all
even.  Every boundary edge is a closed one-cycle on the one-vertex torus.
An odd intersection with one of them proves that the disc boundary has
nonzero class in $H_1(\partial B_t;\mathbb F_2)$.  Its single embedded
boundary circle is thus essential.
\end{proof}

\subsection{Separation from support and coordinate size}

The vector just recovered occupies all $t$ tetrahedra, and its largest
coordinate grows exponentially with $t$.  The argument gives more than
one dense example.

\begin{corollary}[Exact support separation]\label{cor:fib-exact-support}
Every essential standard vertex disc in $B_t$ has quadrilateral support
exactly $t$.  In particular, the minimum such support is $t$, while
$\beta_{\mathrm{LP}}(B_t)=\delta_{\mathrm{LP}}^{\Pi}(B_t)=0$ for every
policy satisfying the deciding-set hypotheses.
\end{corollary}
\begin{proof}
An essential standard vertex disc has a nonzero quadrilateral vector.
It is canonical: a positive vertex-link summand in its standard
coordinates would decompose its extreme ray into nonproportional
nonnegative matching vectors.  Its Euler characteristic is positive,
so the all-anchor propagation argument puts it on
\eqref{eq:fib-recovered-ray}.  Every $W_i$ on that ray is positive.
The equality for $\beta_{\mathrm{LP}}$ was proved above, and the equality
for $\delta_{\mathrm{LP}}^{\Pi}$ follows from
\eqref{eq:deciding-comparison}.
\end{proof}

The previous support theorem supplied an independent quantitative
obstruction: \emph{every} essential standard vertex disc in $B_t$ has
quadrilateral support $a$ satisfying
\begin{equation}
 a\ge\frac{\log_2 F_{t+3}-\log_2 3}{2}.
 \label{eq:fib-support-separation}
\end{equation}
The argument uses the boundary meridional intersection number $F_{t+3}$
and the primitive support-height bound: a disc with support $a$ intersects
any one boundary edge at most $3\cdot4^a$~\cite{priorKernel}.  Every
essential disc in a solid torus has meridional boundary slope, so displaying
one large disc is not the sole basis for the lower bound.  Equation
\eqref{eq:fib-support-separation} excludes all sparse standard vertex
meridians in these triangulations.

The new exact support statement strengthens that earlier
$\Omega(t)$ estimate for this particular family.  The propagation parameter measures
choices among incompatible types; it permits long chains of forced,
nonzero coordinates.  It also permits large integers: the total number of
normal discs in \eqref{eq:fib-standard-vector} is $F_{t+5}-5$, although
each individual coordinate has only $O(t)$ binary digits.  Neither the LP
search nor the ray certificate needs to enumerate those normal discs.
Indeed the complete, uncompressed list of the $7t$ integer coordinates
has $\Theta(t^2)$ binary size: summing the bit lengths of the successive
Fibonacci entries already gives that order.  Once the one-dimensional
cone is known, its coordinates can be reconstructed by $O(t)$ integer
additions with growing $O(t)$-bit operands.  Large geometric multiplicity
and a polynomially sized exact witness coexist here, as required for the
bit-complexity interpretation of the propagation result.

This is an unconditional parameter separation for a specified family,
not a universal bound on $\beta_{\mathrm{LP}}$.  Layered solid tori already
admit specialized direct constructions; we do not claim that they were
previously difficult to recognize.  Their role here is to show rigorously
that the generic mandatory-type mechanism can recover a dense meridian
from the full relaxation, without receiving the correct sector in advance.
